\documentclass[11pt,a4paper]{article}
\usepackage[margin=27mm]{geometry}
\usepackage{amsmath,amssymb,amsthm,hyperref}
\emergencystretch=1em
\newtheorem{lemma}{Lemma}\newtheorem{theorem}[lemma]{Theorem}
\newcommand{\Sz}{\operatorname{Sz}}\newcommand{\gen}[1]{\langle#1\rangle}
\title{Prime subfield residues do not extend to rank four}
\author{Human Compression project}\date{}
\begin{document}\maketitle
\begin{abstract}
We complete the three extension arguments in Section 3.1 of the fixed
higher-rank Suzuki article. None of the prime split-torus rank-three
residues extends to a rank-four family satisfying the source's subgroup
axiom. The proof uses maximal subgroups of an odd dihedral group and
uniqueness of an embedded affine subgroup inside a fixed-size subfield.
In particular, equal abstract orders are never used to identify two
subfield groups. The surviving intermediate candidate in the source's
second case is excluded by the already proved alignment of split tori.
\end{abstract}

\section{Classical and previously proved inputs}
Write $q=2^n$, with $n$ odd, and $G=\Sz(q)$. For a family
$(G_i)$ put $G_J=\bigcap_{j\in J}G_j$. We use exactly the source
condition that for every proper $J$, some unused type gives a maximal
proper inclusion $G_{J\cup\{i\}}<G_J$.
The acquired strict-intersection lemma implies that adding \emph{any}
unused type strictly decreases the intersection. An induced family in
any $G_J$ satisfies the same axiom. No flag-transitivity hypothesis
is imposed.

The classical Suzuki subgroup classification, torus partition and ovoid
model are retained as explicit premises. A nonidentity element in a
split torus has that torus as its centralizer. A full split torus has
normalizer $D_q=C_{q-1}:C_2$. No point stabilizer contains an odd
dihedral group with nontrivial rotation subgroup: in its abelian odd
quotient, an odd rotation has nontrivial image, while an inverting
involution has trivial image, which is impossible.
A proper nonsolvable subgroup is a simple subfield group. A solvable
subgroup is contained in a point stabilizer or one of the torus
normalizers. These are the subgroup-classification consequences already
used in the accepted common notes.

We also reuse the following proved consequences.
\begin{enumerate}
\item Every parabolic in an odd-dihedral family satisfying the subgroup
axiom is maximal in that dihedral group.
\item The full split normalizer of a standard $\Sz(s)$ is
self-normalizing in $G$, and determines exactly one embedded subfield
of a prescribed containing degree. The same-size affine subgroup
$A_s=Z_s:C_{s-1}$ determines exactly one copy of $\Sz(s)$.
\item The proper overgroups of a full split torus in $\Sz(t)$ are
that torus, its split normalizer, the two affine groups at its fixed
points, and the two point stabilizers. A point stabilizer cannot be
a parabolic of a family of rank at least two satisfying the subgroup
axiom.
\item If $t=2^f$ has prime degree $f$ but composite $t-1$, there is
no family of rank at least three satisfying the subgroup axiom.
\item Two full split dihedrals of subfield groups inside a common
split normalizer, sharing an involution, align the two subfields.
Their intersection is the standard group over the intersection field.
\end{enumerate}
These are respectively the accepted cyclic/dihedral lemma, embedded
normalizer and affine uniqueness lemmas, full-torus overgroup lemma,
prime-degree obstruction, and shared-reflection alignment lemma.
Their proofs are imported explicitly; the source's exhaustive geometry
list is not a premise.

\section{One prime-degree maximal subfield}
Suppose $G_0=\Sz(s)$ is maximal in $G$, where $s=2^d\ge8$ and
$p=s-1$ is prime. Then $d$ is an odd prime, and $n=dr$ for an odd
prime $r$. The two primes $d,r$ may coincide. Fix its full split
dihedral $\Delta_s=D_{2p}$ and an affine subgroup $A_s=Z_s:C_p$.

\begin{lemma}\label{overgroups}
A proper subgroup containing $\Delta_s$ and not contained in $G_0$
is contained in the unique ambient split normalizer $D_q$ containing
$\Delta_s$. A proper subgroup containing $A_s$ and not contained in
$G_0$ is contained in the unique point stabilizer fixed by $Z_s$.
\end{lemma}
\begin{proof}
An odd element of order $p$ determines a split torus in $G$.
A simple subfield $\Sz(2^f)$ containing it must have $p\mid2^f-1$:
its nonsplit orders divide $(2^f)^2+1$, which divides $q^2+1$, and
$q-1$ is coprime to $q^2+1$. Since $p=2^d-1$ is prime,
$2$ has order $d$ modulo $p$; it cannot have order one, and $d$ is
prime. Thus $d\mid f$. The only proper divisor $f$ of $dr$ that
is divisible by $d$ is $f=d$. The acquired embedded uniqueness for
$\Delta_s$ or $A_s$ therefore makes such a containing subfield equal
to $G_0$.

A subgroup containing $\Delta_s$ cannot lie in a point stabilizer,
by the odd-dihedral observation, or in a nonsplit normalizer, by torus
coprimality. Its only remaining soluble position is in the unique
$D_q$ fixed by its rotation subgroup. A subgroup containing $A_s$
cannot lie in any torus normalizer: it contains the elementary abelian
$Z_s$ of order $s\ge8$, whereas their Sylow two-subgroups have order
two or four. Its point stabilizer is unique, since each nonidentity
element of $Z_s$ fixes precisely one ovoid point. The subgroup
classification now proves both assertions.
\end{proof}

\begin{lemma}\label{dihedral}
Suppose a rank-four family has $G_{01}=\Delta_s$, and
$G_{013}\cong C_2$. Then
\[
 G_1=D_{2p\ell},\qquad G_{13}=D_{2\ell},
                                                        \tag{1}
\]
where $\ell\ne p$ is an odd prime dividing $(q-1)/p$.
If $G_{012}\cong C_2$, then also
$G_{12}=D_{2\ell}$ and $G_{123}=C_\ell$.
If $G_{012}=C_p$, then $G_{12}$ is either $C_{p\ell}$ or a
dihedral subgroup of order $2p$ distinct from $G_{01}$.
\end{lemma}
\begin{proof}
Strict intersections give $G_1\nleq G_0$. Lemma~\ref{overgroups}
places it inside $D_q$. It contains $\Delta_s$, so is dihedral.
The induced rank-three family in $G_1$ has all parabolics maximal.
Maximality of $\Delta_s$ gives $G_1=D_{2p\ell}$ with $\ell$ prime.
If $\ell=p$, every maximal subgroup of $G_1$ contains its order-$p$
rotation subgroup, contradicting $G_{013}\cong C_2$.
For distinct $p,\ell$, its maximal subgroups are the rotation group
$C_{p\ell}$ and dihedrals of orders $2p$ and $2\ell$.
Only the last type can intersect $G_{01}$ in order two. A dihedral
of order $2\ell$ and $G_{01}$ have a common reflection: their
reflection congruences in the cyclic rotation group have coprime
moduli $p,\ell$. Thus the asserted intersection also has that order.
Two distinct dihedrals of order $2\ell$ in $G_1$ share exactly their
order-$\ell$ rotation group, proving the last intersection claim.
The case of order-$p$ intersection leaves precisely the other two
maximal-subgroup types. Equality of any two displayed parabolics is
excluded by strict intersections.
\end{proof}

\begin{lemma}\label{other}
Under (1), suppose $G_{03}$ is not contained in $G_{01}$. Then
$G_3=\Sz(t)$, where $t=2^r$, $r\ne d$, and $t-1=\ell$ is prime.
\end{lemma}
\begin{proof}
The group $G_3$ contains $G_{13}=D_{2\ell}$. This excludes a
point stabilizer and the nonsplit normalizers. If $G_3$ were soluble,
it would therefore lie in $D_q$, the unique split normalizer containing
this dihedral. But $G_0\cap D_q=\Delta_s$: its odd rotation subgroup
is the full split torus of $G_0$. This would put $G_{03}$ inside
$G_{01}$, a contradiction.

Thus $G_3$ is a proper simple subfield group. Its possible field
exponents are $d$ and $r$, since $n=dr$. The order-$\ell$ rotation
is split in this subfield by the same coprimality argument as above.
The first size is impossible, since $\ell\ne p=s-1$.
Consequently $G_3=\Sz(2^r)$ and $r\ne d$.
Its induced rank-three family satisfies the subgroup axiom. The
prime-degree obstruction implies that $2^r-1$ is prime.
Since it is divisible by $\ell$, it equals $\ell$.
\end{proof}

\section{The three residue types}
\begin{theorem}\label{main}
No rank-four family satisfying the source axiom has a maximal parabolic
$G_0=\Sz(s)$ with $s-1$ prime. In particular none of residues
$R_1,R_2,R_3$ in Section~3.1 of the fixed L110 source extends as
proposed.
\end{theorem}
\begin{proof}
The acquired necessary-pattern theorem gives exactly these three
possibilities in the induced rank-three family of $G_0$. In all cases
we can label $G_{01}=\Delta_s$ and $G_{013}=C_2$.

For $R_1$, we have $G_{02}=A_s$ and $G_{012}=C_2$.
Lemma~\ref{dihedral} gives $G_{12}=D_{2\ell}$.
By strict intersections $G_2\nleq G_0$, so Lemma~\ref{overgroups}
puts $G_2$ in a point stabilizer. But it contains the odd dihedral
$G_{12}$, a contradiction. This proves the first nonextension directly,
without any appeal to an ovoid-circle intersection assumption.

For $R_2$, $G_{02}=A_s$, $G_{03}=D_{2a}$ for an odd prime $a$,
and $G_{012}=C_p$. Again $G_2$ lies in its unique point stabilizer,
so it cannot contain a dihedral group of order $2p$.
Lemma~\ref{dihedral} therefore gives $G_{12}=C_{p\ell}$.
The subgroup $G_{03}$ is not contained in $G_{01}$ by strict
intersections, so Lemma~\ref{other} gives $G_3=\Sz(t)$ with
$t-1=\ell$. The full split dihedrals $G_{01}$ and $G_{13}$ of
$G_0,G_3$ lie in $G_1$ and share $G_{013}=C_2$.
Alignment forces
\[
 G_0\cap G_3\cong\Sz(2^{\gcd(d,r)})=\Sz(2),
\]
which has order twenty. This contradicts $G_{03}=D_{2a}$ for prime
$a$. The source's final $a=5$ candidate is therefore impossible as
well; the extra affine parabolic cannot repair this obstruction.

For $R_3$, $G_{02}=D_{2a}$ and $G_{03}=D_{2b}$, where the
permitted half-orders are primes and at most one is two. All three
intersections $G_{012},G_{013},G_{023}$ have order two, and the
full intersection is trivial. Lemma~\ref{dihedral} gives
\[
 G_{12},G_{13}\cong D_{2\ell},\qquad G_{123}=C_\ell.
\]
Apply Lemma~\ref{other} also with types 2 and 3 interchanged.
Both $G_2,G_3$ are copies of the same $\Sz(t)$, where $t-1=\ell$.
Equal sizes alone do not identify these subgroups.
Instead their common subgroup $G_{23}$ contains the full split torus
$C_\ell$. Within $G_2$, the acquired full-torus overgroup list leaves
$C_\ell$, its normalizer $D_{2\ell}$, the two affine groups, and
the two point stabilizers. The first choice would equal $G_{123}$;
the second would equal $G_{12}$, since the full torus has a unique
normalizer. Both violate strict intersections. A point stabilizer
is excluded as a parabolic of the induced family in $G_2$.
Thus $G_{23}$ is an affine group $A_t=Z_t:C_{t-1}$.
The embedded affine uniqueness theorem now forces its two containing
copies $G_2,G_3$ of $\Sz(t)$ to coincide, the final contradiction.
\end{proof}

\paragraph{Attribution and scope.}
The source is D.~Leemans, \emph{The residually weakly primitive
pre-geometries of the Suzuki simple groups}, Note di Matematica 20(1)
(2000/2001), 1--20, especially printed pp.~10--12. The source's
case division and nonextension conclusions for $R_1,R_3$ are retained;
the proofs here make their overgroup and uniqueness steps explicit.
The $R_2$ obstruction reuses the earlier shared-reflection proof already
applied to the full prescribed lattice of Construction~1.
The general Suzuki structure is classical. No assertion that the other
composite-degree residue lists are exhaustive, no whole-original
completion credit, and no additional search, Lean result or complexity
measurement is made by this component.
\end{document}
